% =====================================================================
%  hookcut - Product & Interface Audit
%  Compiles with pdfLaTeX on Overleaf. No external assets required.
% =====================================================================
\documentclass[11pt,a4paper]{article}

\usepackage[T1]{fontenc}
\usepackage[utf8]{inputenc}
\usepackage[a4paper,margin=25mm,top=24mm,bottom=26mm]{geometry}
\usepackage{libertine}                        % serif body + Biolinum sans
\usepackage[scaled=0.82]{beramono}            % mono
\usepackage{xcolor}
\usepackage{titlesec}
\usepackage{tcolorbox}
\usepackage{booktabs}
\usepackage{tabularx}
\usepackage{enumitem}
\usepackage{fancyhdr}
\usepackage{microtype}
\usepackage[hidelinks]{hyperref}

\tcbuselibrary{breakable,skins}

% ---------- palette ---------------------------------------------------
\definecolor{ink}      {HTML}{16151C}
\definecolor{ink2}     {HTML}{403C4D}
\definecolor{muted}    {HTML}{6D6980}
\definecolor{rule}     {HTML}{DDD9E7}
\definecolor{cardbg}   {HTML}{F6F4FA}
\definecolor{accent}   {HTML}{5B3FD6}
\definecolor{accentbg} {HTML}{EBE6FB}
\definecolor{critcol}  {HTML}{A92C48}
\definecolor{critbg}   {HTML}{FBE9ED}
\definecolor{warncol}  {HTML}{835507}
\definecolor{warnbg}   {HTML}{FBF1DE}
\definecolor{okcol}    {HTML}{2B6A4D}

\color{ink}

% ---------- headings --------------------------------------------------
\titleformat{\section}
  {\sffamily\bfseries\Large\color{ink}}{\thesection}{0.7em}{}
\titleformat{\subsection}
  {\sffamily\bfseries\large\color{ink}}{\thesubsection}{0.6em}{}
\titlespacing*{\section}{0pt}{22pt}{7pt}
\titlespacing*{\subsection}{0pt}{16pt}{5pt}

% ---------- helpers ---------------------------------------------------
\newcommand{\code}[1]{\texttt{\small #1}}
\newcommand{\eyebrow}[1]{%
  {\sffamily\bfseries\footnotesize\color{accent}\MakeUppercase{#1}}\par\vspace{2pt}}

% severity chip
\newcommand{\chip}[3]{%
  \colorbox{#1}{\textcolor{#2}{\sffamily\bfseries\scriptsize\,#3\,}}}

% finding title
\newcommand{\ftitle}[1]{{\sffamily\bfseries\normalsize #1}}

% file reference line
\newcommand{\fref}[1]{\par\vspace{4pt}{\ttfamily\scriptsize\color{muted}#1}}

% finding container - #1 is the accent colour of the left bar
\newtcolorbox{finding}[1]{
  enhanced, breakable,
  colback=cardbg, colframe=#1,
  boxrule=0pt, leftrule=3.5pt,
  arc=0pt, outer arc=0pt,
  left=11pt, right=11pt, top=9pt, bottom=9pt,
  before skip=9pt, after skip=9pt
}

% callout
\newtcolorbox{callout}[1]{
  enhanced, breakable,
  colback=accentbg, colframe=accentbg,
  boxrule=0pt, arc=3pt,
  left=14pt, right=14pt, top=11pt, bottom=11pt,
  before skip=13pt, after skip=13pt,
  title={\textcolor{accent}{\sffamily\bfseries\scriptsize\MakeUppercase{#1}}},
  attach title to upper={\par\vspace{5pt}}
}

% status words in the JTBD table
\newcommand{\stOK}[1]{\textcolor{okcol}{\sffamily\bfseries\small #1}}
\newcommand{\stNo}[1]{\textcolor{critcol}{\sffamily\bfseries\small #1}}
\newcommand{\stPart}[1]{\textcolor{warncol}{\sffamily\bfseries\small #1}}

\setlist[itemize]{leftmargin=1.15em,itemsep=2pt,topsep=4pt}
\setlist[enumerate]{leftmargin=1.4em,itemsep=6pt,topsep=5pt}

% ---------- running head ---------------------------------------------
\pagestyle{fancy}
\fancyhf{}
\renewcommand{\headrulewidth}{0.4pt}
\fancyhead[L]{\sffamily\scriptsize\color{muted}hookcut --- product \& interface audit}
\fancyhead[R]{\sffamily\scriptsize\color{muted}\thepage}

\begin{document}

% =====================================================================
%  TITLE
% =====================================================================
\begin{titlepage}
\vspace*{2.2cm}
\eyebrow{Product \& interface audit}
\vspace{4pt}

{\sffamily\bfseries\fontsize{34}{38}\selectfont
hookcut is two products\\
\textcolor{muted}{wearing one form}\par}

\vspace{16pt}
{\large\color{ink2}
A read of the codebase as it stands: what the tool is built to do, where the
interface stops short of the promise, and what to build first if the next work
is UI/UX.\par}

\vspace{26pt}
\textcolor{rule}{\rule{\linewidth}{0.6pt}}
\vspace{10pt}

{\ttfamily\small\color{muted}
\begin{tabular}{@{}ll@{}}
Repository & Veritron-Innovations/hookcut \\
Commit     & 0c8b27a \\
Surface    & 1{,}395 lines \\
Frontend   & 275 lines across 2 files \\
\end{tabular}\par}

\vfill
{\ttfamily\scriptsize\color{muted}
Read of \code{src/}, \code{backend/} and \code{frontend/} in full.\par}
\end{titlepage}

\tableofcontents
\newpage

% =====================================================================
\section{What hookcut is}

An independent musician or podcaster finishes a track or an episode. To promote
it they need three to five vertical clips for TikTok, Reels and Shorts. hookcut
takes the finished file and produces those clips automatically --- \textbf{cut
from the artist's own audio, not generated}. That distinction is the product's
stated reason to exist, and it is worth holding onto, because nothing in the
current interface communicates it.

Three ideas are doing the work. Whisper produces a timestamped transcript. An
LLM reads that \emph{real} transcript and points at the strongest moments --- it
is told explicitly that it is selecting, not writing. ffmpeg then cuts those
exact timestamps out of the source. The LLM never invents content; it only picks
coordinates.

\begin{callout}{The thesis}
Every other tool in this category generates media. hookcut cuts it. The whole
product rests on the output being genuinely the artist's --- which makes
\emph{provenance} the thing the interface should be showing off, and currently
doesn't show at all.
\end{callout}

% =====================================================================
\section{The pipeline, as built}

Five stages, in a background thread per job, polled every two seconds by the
browser. These stage names are real --- they are the strings the backend writes
and the frontend reads.

\vspace{6pt}
\begin{tabularx}{\linewidth}{@{}l l X@{}}
\toprule
{\sffamily\footnotesize\color{muted}STAGE} &
{\sffamily\footnotesize\color{muted}DOES} &
{\sffamily\footnotesize\color{muted}SOURCE} \\
\midrule
\code{queued}              & Upload accepted          & \code{main.py:104}       \\
\code{transcribing}        & Whisper transcript       & \code{transcribe.py:13}  \\
\code{analyzing}           & Gemini picks moments     & \code{analyze.py:74}     \\
\code{resolving\_cover\_art} & Locate album art       & \code{cover\_art.py:40}  \\
\code{cutting\_clips}      & ffmpeg renders           & \code{cut.py:112}        \\
\bottomrule
\end{tabularx}

\subsection*{Then the pipeline forks on file extension}
\addcontentsline{toc}{subsection}{The fork}

\vspace{4pt}
\noindent
\begin{minipage}[t]{0.48\linewidth}
\begin{tcolorbox}[enhanced,colback=cardbg,colframe=accent,boxrule=0pt,
  toprule=3pt,arc=0pt,left=10pt,right=10pt,top=9pt,bottom=9pt]
{\sffamily\bfseries Audio source --- a song}\par
{\ttfamily\scriptsize\color{muted}.mp3 .wav .m4a .flac .aac .ogg}\par\vspace{6pt}
\begin{itemize}[leftmargin=1em]\small
\item No video exists, so one is \textbf{built}
\item $1080\times1920$ canvas: blurred art fills the frame, sharp art centred
\item Word timings grouped into 7-word lines
\item Karaoke ASS subtitles, word-by-word, with a dim next-line preview
\item Full \code{libx264} re-encode per clip --- \textbf{slow}
\end{itemize}
\end{tcolorbox}
\end{minipage}\hfill
\begin{minipage}[t]{0.48\linewidth}
\begin{tcolorbox}[enhanced,colback=cardbg,colframe=muted,boxrule=0pt,
  toprule=3pt,arc=0pt,left=10pt,right=10pt,top=9pt,bottom=9pt]
{\sffamily\bfseries Video source --- a podcast}\par
{\ttfamily\scriptsize\color{muted}everything else --- .mp4 .mov}\par\vspace{6pt}
\begin{itemize}[leftmargin=1em]\small
\item Stream copy, \code{-c copy} --- near-instant
\item Keeps source aspect ratio, so typically \textbf{16:9}
\item No vertical reframe
\item No lyrics, no subtitles, no overlay
\item Cover art step skipped entirely
\end{itemize}
\end{tcolorbox}
\end{minipage}

\subsection{That fork is the most important fact about this UI}

The two branches produce fundamentally different artefacts. The song path
returns a finished vertical video with synced lyrics burned in --- genuinely
close to postable. The podcast path returns a horizontal excerpt of the original
file, which is a long way from postable: it still needs reframing, captions and
text before it can go anywhere.

\textbf{The interface treats both as the same thing.} One form, one set of
controls, one card layout, one aspect ratio. A musician and a podcaster see an
identical screen and get wildly different value from it. Any serious UI/UX work
has to decide, early, whether hookcut is one product with a branch or two
products sharing a pipeline --- and the honest answer today is the latter.

% =====================================================================
\section{What the project should accomplish}

Doing this by hand is roughly seven steps and two to four hours per release:
listen back and pick moments, cut them, build a vertical background, transcribe
and sync the lyrics, add a hook text overlay, write per-platform captions,
export each clip. hookcut's goal is to collapse that into about five minutes.

The test is simple and unforgiving: \textbf{can the artist post the clip without
opening another app?} Any step that sends them to CapCut is the product failing
its own promise. Measured that way:

\vspace{8pt}
\begin{tabularx}{\linewidth}{@{}p{3.1cm} l X@{}}
\toprule
{\sffamily\footnotesize\color{muted}MANUAL STEP} &
{\sffamily\footnotesize\color{muted}TODAY} &
{\sffamily\footnotesize\color{muted}WHERE IT STANDS} \\
\midrule
Pick the moments &
\stOK{Automated} &
\small The LLM chooses, and is prompted to always include the chorus. But the
choice is final --- nothing can be adjusted. \\[4pt]

Cut the audio &
\stOK{Automated} &
\small Exact ffmpeg cuts at the chosen timestamps. \\[4pt]

Vertical background &
\stOK{Automated} &
\small Songs only. One fixed look --- blurred art behind sharp art. No
alternatives. \\[4pt]

Sync the lyrics &
\stOK{Automated} &
\small Songs only. Karaoke-timed with a next-line preview. Whisper's word timing
on sung audio is approximate, and there is no way to correct it. \\[4pt]

Hook text overlay &
\stNo{Suggested only} &
\small Three variants written per clip, then printed as bullet points. Nothing
renders them. \textbf{The artist goes to CapCut.} \\[4pt]

Write captions &
\stPart{Not usable} &
\small TikTok and IG captions are produced, then displayed as plain text with no
copy control. \\[4pt]

Export and post &
\stPart{Download only} &
\small A download link per clip. No batch download, no share target. \\
\bottomrule
\end{tabularx}

\vspace{10pt}
Four of seven steps land. Three push the artist back out of the app --- and the
one that matters most, the text overlay, is the single thing standing between
``here are some clips'' and ``here is a post.''

% =====================================================================
\section{The promise gap}

\begin{finding}{critcol}
\chip{critbg}{critcol}{P0}\quad\ftitle{Text overlays are written but never rendered}
\par\vspace{5pt}
\code{analyze.py} returns three overlay variants per clip, each aimed at a
different trigger --- relatability, curiosity, contrarian. The frontend prints
them as bullets. No code path anywhere consumes them. \code{render\_video.py}
builds an ASS subtitle file already, so the machinery to burn text into the
frame exists and is running; the overlay simply is not wired into it.
\fref{analyze.py:65 $\rightarrow$ page.tsx:210--214 \textbullet{} dead end}
\end{finding}

\begin{finding}{critcol}
\chip{critbg}{critcol}{P0}\quad\ftitle{Captions cannot be copied}
\par\vspace{5pt}
Both captions render as bare paragraphs. The user selects text with a mouse and
copies it by hand, twice per clip, five clips per run. A copy button is a
trivial addition and removes ten manual interactions per job.
\fref{page.tsx:218, page.tsx:223}
\end{finding}

\begin{finding}{warncol}
\chip{warnbg}{warncol}{P1}\quad\ftitle{No batch export}
\par\vspace{5pt}
Five clips means five separate download clicks, each landing a file named
\code{clip\_2\_the\_confession.mp4}. No zip, no naming the artist would
recognise.
\fref{page.tsx:226 \textbullet{} cut.py:136}
\end{finding}

% =====================================================================
\section{Control and trust}

\begin{finding}{critcol}
\chip{critbg}{critcol}{P0}\quad\ftitle{One shot, no preview, no recourse}
\par\vspace{5pt}
Upload, wait several minutes, receive five finished clips. There is no point at
which the artist sees the transcript, sees which moments were chosen, or can say
``not that one.'' If the LLM picks badly, the only available action is
\code{reset()} --- which clears everything and demands a fresh upload and a
fresh transcription of the same file.
\par\vspace{5pt}
The pipeline has a natural seam the UI ignores: transcribe and analyse are
comparatively fast and cheap; rendering is slow and expensive. Splitting there
turns a gamble into a decision.
\fref{page.tsx:78--84 \textbullet{} main.py:53--101}
\end{finding}

\begin{finding}{critcol}
\chip{critbg}{critcol}{P0}\quad\ftitle{The evidence is fetched, then thrown away}
\par\vspace{5pt}
Every concept carries \code{source\_text} --- the actual transcript line the
moment was chosen for. The frontend's own type declares it. \code{ClipCard}
renders the angle name, the timestamps, the overlays and both captions, and
silently drops \code{source\_text}.
\par\vspace{5pt}
This is the provenance signal. It is the one field that proves the clip came
from the artist's real material rather than being invented --- the product's
entire differentiator --- and it is the one field the card does not show.
\fref{page.tsx:11 declared \textbullet{} page.tsx:194--231 never used}
\end{finding}

\begin{finding}{warncol}
\chip{warnbg}{warncol}{P1}\quad\ftitle{Cut points are fixed, and clips are all-or-nothing}
\par\vspace{5pt}
The LLM's \code{mm:ss} range goes straight to ffmpeg. A clip that starts half a
word early cannot be nudged, and a single bad clip cannot be re-rolled without
re-running all five.
\fref{cut.py:135--158}
\end{finding}

% =====================================================================
\section{Progress theatre}

\begin{finding}{critcol}
\chip{critbg}{critcol}{P0}\quad\ftitle{The progress bar is a static rectangle}
\par\vspace{5pt}
It is hardcoded to \code{width: "60\%"} and given
\code{animation: "pulse 1.5s infinite"}. There is no \code{pulse} keyframe
anywhere in the repository --- and no CSS file in which to define one. So it
renders as a motionless bar, stuck at 60\%, for the entire multi-minute run. It
reads as a hung process.
\fref{page.tsx:164--166 \textbullet{} no @keyframes in repo}
\end{finding}

\begin{finding}{critcol}
\chip{critbg}{critcol}{P0}\quad\ftitle{Five known stages, none of them counted}
\par\vspace{5pt}
The backend reports exactly which stage it is in and the frontend has a label
for each. Nothing tells the user it is stage three of five, how long they have
been waiting, or that transcription alone can take minutes. During
\code{cutting\_clips} --- the longest stage --- there is no per-clip count
either, though the loop knows the index.
\fref{page.tsx:24--32 \textbullet{} cut.py:135}
\end{finding}

\begin{finding}{warncol}
\chip{warnbg}{warncol}{P1}\quad\ftitle{Uploads are silent}
\par\vspace{5pt}
An hour-long podcast is a large file. The submit handler uses \code{fetch},
which reports no upload progress, so the button says ``Uploading\ldots'' and
nothing else happens until the whole file has landed. Upload progress needs
\code{XMLHttpRequest}.
\fref{page.tsx:70}
\end{finding}

% =====================================================================
\section{Three ways it hangs without saying so}

These are ordinary bugs, but they are listed here because each one presents to
the user as a frozen interface with no error --- the worst possible failure mode
for a process that is \emph{supposed} to take minutes.

\begin{finding}{critcol}
\chip{critbg}{critcol}{P0}\quad\ftitle{A dead backend locks the button forever}
\par\vspace{5pt}
\code{handleSubmit} has no error handling. If the fetch rejects --- backend not
running, connection refused --- the line that clears \code{submitting} never
executes. The button stays disabled reading ``Uploading\ldots'' permanently,
with no error shown. Only a page reload recovers it.
\fref{page.tsx:57--76 \textbullet{} setSubmitting(false) at :72 unreachable on throw}
\end{finding}

\begin{finding}{critcol}
\chip{critbg}{critcol}{P0}\quad\ftitle{A server restart produces an infinite spinner}
\par\vspace{5pt}
Jobs live in an in-memory dict. When the backend restarts, that dict is empty,
and the poll endpoint answers with \code{\{"error": "job not found"\}} --- an
object with \textbf{no \code{stage} field}. The frontend stores it, finds
\code{stage} is neither \code{done} nor \code{error}, and so keeps rendering the
in-progress view and keeps polling. Forever.
\par\vspace{5pt}
The README tells developers to run uvicorn with \code{-{}-reload}, so in normal
development this fires on every file save.
\fref{main.py:50, main.py:144 $\rightarrow$ page.tsx:48--53}
\end{finding}

\begin{finding}{warncol}
\chip{warnbg}{warncol}{P1}\quad\ftitle{Failed polls throw into the void}
\par\vspace{5pt}
The polling callback is unguarded. A single failed request throws inside the
interval, which keeps firing regardless, accumulating rejections that never
reach the user.
\fref{page.tsx:46--55}
\end{finding}

% =====================================================================
\section{One form, two paths}

\begin{finding}{warncol}
\chip{warnbg}{warncol}{P1}\quad\ftitle{Podcast clips are forced into a vertical box}
\par\vspace{5pt}
Every clip card hardcodes \code{aspectRatio: "9/16"}. Song clips genuinely are
$1080\times1920$, so they fit. Podcast clips are stream copies at the source
aspect --- usually 16:9 --- and get letterboxed into a tall black frame with a
slender strip of video in the middle.
\fref{page.tsx:200 vs cut.py:75}
\end{finding}

\begin{finding}{warncol}
\chip{warnbg}{warncol}{P1}\quad\ftitle{Controls that explain when they do not work}
\par\vspace{5pt}
The lyrics checkbox is labelled ``Burn in karaoke-style lyrics
\emph{(audio-only sources)}'' and is always shown --- including when the user
has just selected an mp4, where it does nothing. The cover-art upload has the
same problem. The file type is known the moment it is chosen; the form could
simply become the right form.
\fref{page.tsx:106--114, page.tsx:150--153}
\end{finding}

\begin{finding}{accent}
\chip{accentbg}{accent}{P2}\quad\ftitle{Invisible defaults}
\par\vspace{5pt}
Genre and mood are optional free-text that silently fall back to ``music'' and
``moody''. A user leaving them blank has no idea a default was applied or what
it was. Clip count is a number stepper for a range of 1--10, where a segmented
control would take one click instead of several.
\fref{page.tsx:64--65, page.tsx:138--148}
\end{finding}

% =====================================================================
\section{There is no styling foundation}

\begin{finding}{warncol}
\chip{warnbg}{warncol}{P1}\quad\ftitle{Zero CSS files, and it constrains everything}
\par\vspace{5pt}
The repository contains no stylesheet of any kind --- no CSS file, no Tailwind,
no CSS modules, no \code{next.config.js}, no PostCSS. All styling is inline
React style objects. This is not a taste question; inline styles structurally
cannot express several things the interface needs:
\begin{itemize}
\item \textbf{No pseudo-classes.} No hover, focus, active or disabled styling.
      Buttons do not react to the cursor at all.
\item \textbf{No media queries.} The layout is a fixed 720px column. The results
      grid is the only responsive element in the app --- on a phone, the form
      simply does not adapt.
\item \textbf{No keyframes.} Which is why the progress animation is dead.
\item \textbf{No focus-visible.} Keyboard navigation is invisible.
\item \textbf{No tokens.} The same eight hex values are retyped throughout, and
      the dark theme is hardcoded on the body element.
\end{itemize}
\vspace{3pt}
Everything else in the P1 tier is cheaper to fix once this is fixed first.
\fref{page.tsx:233--261 \textbullet{} layout.tsx:9 \textbullet{} no *.css in repo}
\end{finding}

% =====================================================================
\section{Accessibility}

\begin{finding}{warncol}
\chip{warnbg}{warncol}{P1}\quad\ftitle{Stage changes are never announced}
\par\vspace{5pt}
The status region has no \code{aria-live}, so a screen-reader user gets no
notification when the job moves from transcribing to analysing to done. On a
process this long, that is the entire feedback channel missing.
\fref{page.tsx:161--168}
\end{finding}

\begin{finding}{warncol}
\chip{warnbg}{warncol}{P1}\quad\ftitle{The progress bar is not a progress bar}
\par\vspace{5pt}
It is two nested \code{div}s with no \code{role="progressbar"} and no value
attributes --- invisible to assistive technology.
\fref{page.tsx:164--166}
\end{finding}

\begin{finding}{accent}
\chip{accentbg}{accent}{P2}\quad\ftitle{Contrast, at least, is fine}
\par\vspace{5pt}
Worth recording as a pass: the muted grey \code{\#9a9aa5} on both the page
ground and the card surface clears WCAG AA comfortably. The palette is sound ---
it is the interaction states that are missing, not the colours.
\fref{page.tsx:255--261}
\end{finding}

% =====================================================================
\section{State and continuity}

\begin{finding}{warncol}
\chip{warnbg}{warncol}{P1}\quad\ftitle{Refreshing the page abandons the job}
\par\vspace{5pt}
The job id lives only in component state --- not in the URL, not in storage. A
refresh, an accidental navigation, or a closed tab loses all access to a job
that is still running happily on the server. Putting the id in the URL makes the
run survivable and shareable, and costs almost nothing.
\fref{page.tsx:41}
\end{finding}

\begin{finding}{accent}
\chip{accentbg}{accent}{P2}\quad\ftitle{No history, no cancel}
\par\vspace{5pt}
Past runs are unreachable even though their files are still on disk under
\code{backend/jobs/}. A running job cannot be cancelled --- a mistaken upload
occupies the machine until it finishes.
\fref{main.py:43--50}
\end{finding}

% =====================================================================
\section{What I would build first}

Ordered by how much each one moves the product toward ``the artist posts without
opening another app.'' The first three are the ones that change what hookcut
\emph{is}; the rest make it feel finished.

\begin{enumerate}
\item \textbf{Split the flow at the review seam.}\\
      \small\color{ink2}
      Transcribe and analyse, then stop and show the artist what was chosen ---
      transcript, moments, source lines. They approve, drop or adjust, and only
      then does rendering start. Turns the current gamble into a decision, and
      stops wasting minutes of encoding on clips nobody wanted.

\item \textbf{Burn the chosen overlay into the clip.}\\
      \small\color{ink2}
      Let the artist pick one of the three suggested hooks; render it into the
      frame. \code{render\_video.py} already writes and burns ASS subtitles, so
      this is a new style and a headline event, not new infrastructure. This
      single change is what closes the CapCut gap.

\item \textbf{Make waiting honest, and make failure visible.}\\
      \small\color{ink2}
      Real stage-of-five progress with elapsed time and a per-clip count, plus
      the three hang bugs fixed: guard the submit, guard the poll, and treat a
      \code{stage}-less response as an error rather than as ``still working.''

\item \textbf{Show the source line on every card.}\\
      \small\color{ink2}
      \code{source\_text} is already in the payload. Displaying it is a few
      lines and it is the product's whole argument --- this really is your
      material, here is the line it came from, here is where it sits in the
      track.

\item \textbf{Give the frontend a stylesheet.}\\
      \small\color{ink2}
      Tailwind or a tokenised \code{globals.css} --- either works. This unblocks
      hover and focus states, real responsiveness, working animation and
      keyboard visibility all at once, and every subsequent UI change gets
      cheaper.

\item \textbf{Let the form follow the file.}\\
      \small\color{ink2}
      On selection, branch: a song gets cover art and lyrics controls and a 9:16
      preview; a podcast gets neither and a correct aspect ratio. Set the right
      expectation about what comes out, per path.

\item \textbf{Copy buttons, batch download, job id in the URL.}\\
      \small\color{ink2}
      Small, cheap, immediately felt --- the difference between a demo and a
      tool someone uses weekly.
\end{enumerate}

\begin{callout}{If only one thing gets built}
Items 1 and 2 together. Review-then-render plus a burned-in hook is the
difference between hookcut producing \emph{clips} and hookcut producing
\emph{posts} --- and everything else on this list is polish on top of whichever
of those two the product turns out to be.
\end{callout}

% =====================================================================
\section{Notes for later}

Not UI/UX work, but they would bite during it, so they are recorded here rather
than lost.

\begin{itemize}
\item \textbf{Confirm the model id resolves.} \code{analyze.py:102} calls
      \code{gemini-3.6-flash}. If that name does not resolve on the account
      being used, every job dies at the analysing stage --- worth a single check
      before debugging the frontend, because the symptom would look like a UI
      bug.
\item \textbf{\code{/clips} serves the whole jobs directory.} The static mount
      exposes not just the rendered clips but the original uploaded file and the
      transcript, to anyone holding the job UUID. Harmless on localhost; a
      blocker before any deployment.
\item \textbf{Uploaded filenames go straight into a path.} \code{main.py:117}
      joins the client-supplied filename without sanitising it.
\item \textbf{The backend is explicitly single-user.} Its own docstring says to
      swap the thread-per-request model for a real queue before multi-user use
      --- worth honouring that note before any UI implies concurrency.
\item \textbf{Local setup is missing two things.} Python is not installed on the
      audited machine, only the Microsoft Store alias stub, and ffmpeg is not
      installed at all. Node is fine. The frontend can be run and styled without
      either, but nothing can be rendered end-to-end until both are in place.
\end{itemize}

\end{document}
